\chapter*{Control del documento}
\addcontentsline{toc}{chapter}{Control del documento}

\begin{tabularx}{\textwidth}{>{\bfseries}p{4.1cm}Y}
Título & Proyecto Final: Fase 1 - Diseño, Subneteo e Implementación Base en la Nube \\
Organización propuesta & RedNova GT \\
Curso & CC3067 Redes, sección 10 \\
Equipo & Grupo 1 \\
Versión & 1.3 \\
Fecha objetivo & 19 de octubre de 2026 \\
Clasificación & Documento académico \\
Repositorio & \href{https://github.com/DanielBarillasM/Proyecto-final_Grupo-1_Redes_Sec-10.git}{\texttt{GitHub: Proyecto Final - Grupo 1}} \\
\end{tabularx}

\section*{Historial de revisión}

\begin{tabularx}{\textwidth}{C{1.65cm} C{2.4cm} Y Y}
\toprule
\textbf{Versión} & \textbf{Fecha} & \textbf{Descripción} & \textbf{Responsable} \\
\midrule
0.1 & 5/10/2026 & Definición de alcance, supuestos y arquitectura preliminar. & Grupo 1 \\
0.7 & 5/10/2026 & Incorporación de cálculos, diagramas, Terraform y controles de seguridad. & Grupo 1 \\
1.0 & 5/10/2026 & Compilación, validación técnica y revisión visual del entregable. & Grupo 1 \\
1.1 & 19/10/2026 & Ajustes editoriales solicitados, nueva captura PowerShell y auditoría final. & Grupo 1 \\
1.2 & 20/10/2026 & Evidencia de identidad y planificación en AWS, enlace al repositorio y actualización de pendientes. & Grupo 1 \\
1.3 & 20/10/2026 & Nota de ejecución final para controlar costos y actualización del riesgo de credenciales temporales. & Grupo 1 \\
\bottomrule
\end{tabularx}

\section*{Integrantes}

\begin{tabularx}{\textwidth}{Y C{3cm}}
\toprule
\textbf{Nombre} & \textbf{Carné} \\
\midrule
Pablo Daniel Barillas Moreno & 22193 \\
Andrés Rafael Chivalán & 21534 \\
Roberto Emiliano Otoniel Camposeco Torres & 23968 \\
\bottomrule
\end{tabularx}

\section*{Declaración de trazabilidad}

El informe distingue tres clases de afirmaciones: diseño propuesto, validación local reproducible y evidencia obtenida en servicios externos. Los resultados de herramientas se conservan en la carpeta \code{evidence/}; las credenciales, llaves privadas y estados de Terraform están excluidos del control de versiones. Las capturas añadidas acreditan autenticación y planificación contra AWS, pero no se presentan como evidencia de un \code{apply}, de recursos creados ni de conectividad entre instancias.
